% Fichier complété par tools/complete_missing_corrections.py.
% Source : ELEC/ELEC-01/537_CircuitElec/537_CircuitElec.tex
% Statut : rédigé (1 question(s)).
\begin{corrigebox}[Corrigé rédigé]
\CorrigeQuestion{1}
On réduit le réseau depuis l'extrémité opposée aux bornes. Deux
            résistances en série donnent $R_s=R_a+R_b$ et deux résistances en parallèle
            $R_p=R_aR_b/(R_a+R_b)$. Pour le réseau en échelle, on définit récursivement
            $Z_{n}=R_{v,n}$ puis
            \[Z_k=R_{v,k}\parallel(R_{h,k+1}+Z_{k+1}),\qquad
              R_{eq}=R_{h,1}+Z_1.\]
            Pour le pont symétrique, les nœuds équipotentiels peuvent être fusionnés avant
            la réduction.
\end{corrigebox}
